\section{Introduction}

The topical collection \emph{Geometry of Classical and Quantum Space-Times} foregrounds a recurring theme in gravitational physics: geometry is not merely a stage on which physics happens, but a structure that may itself be reconstructed, constrained, or even \emph{defined} by operational and information-theoretic considerations. This review is written in that spirit, and it is intentionally \emph{geometry-first}: we begin from well-posed geometric objects and dynamical systems and then introduce an explicit \emph{readout protocol} that turns those objects into finite-resolution data streams. The central organizing idea is that ``spacetime'' can be approached as an emergent interface---a stable quotient of a richer geometry under a constrained observation map---rather than as a primitive background.

\subsection{Scope and guiding separation of layers}

Throughout, we maintain a strict separation between:
\begin{enumerate}[label=(\roman*),leftmargin=*,itemsep=2pt]
  \item \textbf{Object layer} (geometry/dynamics): superspace embeddings, lattices, torus translations, cut-and-project schemes, model sets, and their invariant measures;
  \item \textbf{Protocol layer} (readout): finite-resolution windowing, coarse-grainings, symbolic codings, canonicalization/stabilization maps, and the induced fibers of unresolved degrees of freedom;
  \item \textbf{Interpretation layer} (physics): how the protocol-level objects can be mapped to familiar spacetime notions (time, causality, observables), and where such mappings should \emph{not} be over-claimed.
\end{enumerate}
This separation is essential for a review that aims to be useful across classical and quantum space-time contexts: it prevents category errors (e.g.\ treating a protocol artifact as an ontological claim) while keeping the mathematical backbone auditable.

\subsection{What is covered}

The review is centered on a concrete chain of constructions:
\begin{itemize}[leftmargin=*,itemsep=2pt]
  \item \textbf{Superspace model sets}: a six-dimensional periodic structure and a choice of physical/internal decomposition leading to a cut-and-project model set; acceptance domains (``windows'') and atomic surfaces provide the geometry-to-pattern interface.
  \item \textbf{Finite-resolution readout}: the observer is modeled by parameters (orientation, internal phase, window shape/width, and resolution kernel). Readout maps continuous internal degrees of freedom to discrete symbols.
  \item \textbf{Stabilization and fibers}: canonicalization maps (Fold-type) compress raw readout words into stable languages (e.g.\ golden-mean constraints), inducing preimage fibers that encode unresolved degrees of freedom.
  \item \textbf{Information-theoretic time}: ``time'' is defined protocol-wise by refinement of posterior cylinder sets; under ergodic hypotheses, the contraction rate is controlled by entropy rates (Shannon--McMillan--Breiman / Kolmogorov--Sinai).
  \item \textbf{Causality interface}: we discuss a conservative interface to Lorentzian causality, including a lightcone-boundary readout viewpoint formulated at the protocol layer.
\end{itemize}

\subsection{Main contributions of this review}

As a review article, the primary contribution is synthesis. Nevertheless, we aim to contribute \emph{structure} rather than mere exposition. Concretely:
\begin{enumerate}[label=\textbf{C\arabic*.},leftmargin=*,itemsep=2pt]
  \item \textbf{A unified dependency chain}: we present a reusable minimal interface linking superspace model sets, finite-resolution readout, stabilization/canonical forms, fibers, and entropy-rate time, with each link stated as a definition/assumption/proposition template.
  \item \textbf{A protocol-level definition of time}: we emphasize an ``entropy-rate time'' viewpoint that does not presuppose a continuum time parameter at the object layer, but arises from posterior refinement under readout.
  \item \textbf{A cautious causality bridge}: we propose a clean interface between readout-induced partial orders/constraints and Lorentzian-causal language, clarifying what is definitional, what is conjectural, and what is testable.
  \item \textbf{Falsifiable fingerprints}: we collect observable statistical signatures that are natural from the superspace/readout viewpoint and are, in principle, accessible to numerical tests or laboratory analogues.
\end{enumerate}

\subsection{Roadmap}

Section~\ref{sec:superspace} introduces superspace model sets as geometric primitives. Section~\ref{sec:readout} formalizes finite-resolution readout as an operator/protocol acting on internal degrees of freedom. Section~\ref{sec:stabilization} develops stabilization maps and the associated fiber structures. Section~\ref{sec:info-time} defines information-theoretic time via refinement and entropy rates. Section~\ref{sec:causality} discusses causality interfaces. Section~\ref{sec:falsifiability} lists falsifiable fingerprints. We close with limitations and open problems in Section~\ref{sec:discussion}.

\subsection{Positioning within spacetime reconstruction programs}\label{sec:intro:positioning}

Spacetime reconstruction appears across a broad range of classical and quantum gravity approaches. To help GRG readers calibrate the scope of this review, we briefly situate the present ``superspace readout'' viewpoint relative to several established themes:
\begin{itemize}[leftmargin=*,itemsep=2pt]
  \item \textbf{Causal-set and causal-structure-first approaches}: in causal set theory, causal order and discreteness are taken as fundamental, with continuum Lorentzian geometry expected to emerge in appropriate limits \cite{Bombelli1987,Sorkin2003}. Our review does not postulate a causal order at the object layer; instead, we define a protocol precedence relation induced by refinement (Section~\ref{sec:causality}) and discuss what would be needed to relate it to Lorentzian causal orders.
  \item \textbf{Quantum geometry via operator algebras}: approaches such as loop quantum gravity emphasize that geometric information is encoded in (noncommuting) operators and their spectra \cite{AshtekarLewandowski2004,Thiemann2007}. Our ``readout as operator/protocol'' language resonates with this viewpoint, but we remain agnostic about quantization: the noncommutativity we discuss can be treated as a protocol constraint even in a purely classical object layer.
  \item \textbf{Spectral and information-theoretic reconstructions}: a recurring idea is that geometric or causal information can be inferred from spectral data, entropies, or information-flow constraints. The entropy-rate time defined here (Section~\ref{sec:info-time}) belongs to this family, but it is grounded in the SMB/KS theorem chain and tied explicitly to a readout protocol.
  \item \textbf{Superspace and higher-dimensional embeddings}: in quasicrystal theory, higher-dimensional periodic embeddings provide a mature technology for generating aperiodic patterns with controlled harmonic structure \cite{BaakeGrimm2013}. We import this technology as a geometric backbone and then ask how protocol-limited observation turns it into ``spacetime-like'' notions.
\end{itemize}

The practical message is that this article is not a proposal of a new quantum gravity model. Rather, it is a structured review of a reusable interface: \emph{superspace geometry + finite-resolution readout} as a disciplined generator of symbolic data streams, stable types, fiber structures, and entropy-controlled arrows. This interface can be used as a building block inside different reconstruction programs, provided its assumptions and limits are kept explicit.

\subsubsection{Why a superspace backbone is useful for GR/QG-facing discussions}

From a GR/QG standpoint, one might ask why a tool originating in quasicrystal geometry should be relevant at all. The reason is not that superspace ``is spacetime,'' but that superspace model sets offer a rare combination of properties that are attractive for reconstruction:
\begin{itemize}[leftmargin=*,itemsep=2pt]
  \item \textbf{Exact periodicity at the object layer}: the higher-dimensional description is periodic and therefore amenable to harmonic analysis, ergodic theorems, and explicit dual modules.
  \item \textbf{Controlled emergence of aperiodicity}: aperiodicity is generated by a projection/cut mechanism that is explicit and parameterized, rather than by ad hoc randomness. This supports auditable links between geometry and statistics.
  \item \textbf{Natural separation into visible and hidden degrees of freedom}: the physical/internal split is built in, making it straightforward to define what is ``accessible'' only through readout.
\end{itemize}
These features align well with the ethos of many geometry-first programs: one seeks structures whose macroscopic behavior is constrained by geometry and symmetry, while the role of observation is made explicit and does not silently assume a continuum limit.

\subsubsection{A note on the Lewandowski context}

The memorial context of the topical collection naturally emphasizes rigorous geometric structures and their quantization. While this review does not attempt to connect directly to specific constructions in loop quantum gravity, it shares the guiding idea that one should avoid background-dependent crutches and should separate kinematics (what the objects are) from dynamics (how they evolve) and from interpretation (how they map to physical observables) \cite{AshtekarLewandowski2004}. In that spirit, the present article focuses on kinematical geometry and on protocol-defined observables and arrows; dynamics beyond the scan/readout level is left as an explicitly identified open problem.
